% =========================
% E) Evaluation Metrics and Comparison Targets
% =========================
\subsection{Evaluation Metrics and Comparison Targets}
\label{subsec:metrics}

\paragraph{Performance traces.}
For run \(r\), let \(y_t^{(r)}\) denote the observed objective value at evaluation \(t\)
(or the objective vector \(\mathbf{y}_t^{(r)}\) in the multiobjective case).
For minimization, the incumbent best-so-far value is
\begin{equation}
    f_{\min}^{(r)}(t) = \min_{1 \le i \le t} y_i^{(r)}.
\end{equation}
In the multiobjective setting, define the set of observed vectors
\[
    \mathcal{S}^{(r)}(t) = \{\mathbf{y}_1^{(r)},\dots,\mathbf{y}_t^{(r)}\},
\]
and let
\[
    \mathcal{N}^{(r)}(t) = \mathrm{ND}\!\big(\mathcal{S}^{(r)}(t)\big)
\]
denote its nondominated subset.

\paragraph{Oracle targets (empirical reference).}
When the true optimum / Pareto front is unknown, we construct an \emph{empirical oracle reference} at the largest tested budget \(\BmaxSymbol\):
\begin{itemize}
    \item \textbf{Single-objective oracle:}
    \[
        f^\star = \min_{\text{all methods},\ \text{all runs}} f_{\min}^{(r)}(\BmaxSymbol).
    \]
    \item \textbf{Multiobjective oracle front:}
    \[
        \mathcal{N}^\star =
        \mathrm{ND}\!\left(
        \bigcup_{\text{all methods},\ \text{all runs}} \mathcal{S}^{(r)}(\BmaxSymbol)
        \right).
    \]
\end{itemize}
This oracle is a benchmarking reference (best observed value/front), not a claim of global optimality.

% ---------------------------------------------------------
\subsubsection{Single-objective metrics (for \(m=1\))}

\paragraph{Fixed-budget regret (fb).}
At budget \(B\), regret is the gap to the best observed oracle value:
\begin{equation}
    \mathrm{Regret}^{(r)}(B) = f_{\min}^{(r)}(B) - f^\star.
\end{equation}

\paragraph{Fixed-target best (ftb): evaluations to reach oracle-best.}
The evaluation count required to match the oracle-best value is
\begin{equation}
    T_{\mathrm{ftb}}^{(r)} =
    \min\left\{t \in \{1,\dots,\BmaxSymbol\} : f_{\min}^{(r)}(t) \le f^\star \right\},
\end{equation}
with \(T_{\mathrm{ftb}}^{(r)}=\infty\) if the target is never reached within \(\BmaxSymbol\) evaluations.

\paragraph{Fixed-target close (ftc): evaluations to reach near-oracle.}
We define a near-oracle threshold
\[
    f^{\dagger} = f^\star + \delta,
    \qquad
    \delta = \NearOracleEta |f^\star|
\]
unless otherwise specified.
The corresponding time-to-threshold is
\begin{equation}
    T_{\mathrm{ftc}}^{(r)} =
    \min\left\{t : f_{\min}^{(r)}(t) \le f^{\dagger} \right\},
\end{equation}
with \(\infty\) if never reached.

% ---------------------------------------------------------
\subsubsection{Multiobjective metrics (for \(m>1\))}

\paragraph{Dominance and recovery criteria.}
We evaluate multiobjective recovery in normalized objective space (per \ObjectiveNormalizationPolicy{}).
A point is considered ``recovered'' if it matches an oracle point within tolerance \(\tau\):
\[
    \|\tilde{\mathbf{y}} - \tilde{\mathbf{y}}^\star\|_{\infty} \le \tau,
    \qquad
    \tau=\RecoveryTolerance.
\]

\paragraph{Fixed-budget hypervolume (fbp-HV).}
After monotone normalization, let \(\tilde{\mathcal{N}}^{(r)}(B)\) denote the nondominated set at budget \(B\) in normalized space.
We compute
\begin{equation}
    \mathrm{HV}^{(r)}(B)
    = \mathrm{HV}\!\big(\tilde{\mathcal{N}}^{(r)}(B); \mathbf{r}\big),
\end{equation}
with reference point \(\mathbf{r}\) specified by the normalization policy.
We also report the hypervolume gap
\[
    \Delta \mathrm{HV}^{(r)}(B) = \mathrm{HV}^\star - \mathrm{HV}^{(r)}(B),
\]
where \(\mathrm{HV}^\star\) is the hypervolume of the oracle set \(\tilde{\mathcal{N}}^\star\).

\paragraph{Fixed-budget additive \(\epsilon\)-indicator (fbp-\(\epsilon\)).}
We compute the additive \(\epsilon\)-indicator relative to the oracle front:
\begin{equation}
    I_{\epsilon}^{(r)}(B) =
    \inf\left\{\epsilon :
    \forall \mathbf{u}\in\tilde{\mathcal{N}}^\star,\
    \exists \mathbf{v}\in\tilde{\mathcal{N}}^{(r)}(B)
    \ \text{s.t.}\ v_i \le u_i + \epsilon\ \forall i
    \right\}.
\end{equation}

\paragraph{Fixed-target one (fto): evaluations to recover at least one oracle point.}
We define
\begin{equation}
    T_{\mathrm{fto}}^{(r)} =
    \min\left\{t :
    \exists \mathbf{y}\in \mathcal{N}^{(r)}(t)\ \text{recovering some}\ \mathbf{y}^\star\in\mathcal{N}^\star
    \right\}.
\end{equation}

\paragraph{Fixed-target all (fta): evaluations to recover the oracle front (approx.).}
Exact recovery may be unrealistic under continuous/noisy objectives.
We therefore use a coverage criterion: recover at least fraction \(q=\CoverageFraction\) of oracle points
(under the recovery tolerance \(\tau\)). We report \(q\) and \(\tau\) explicitly.

\paragraph{Primary vs.\ secondary metrics.}
\textbf{Primary:} fixed-budget regret (single objective) or hypervolume gap \(\Delta\mathrm{HV}\) (multiobjective).
\textbf{Secondary:} fixed-target times (\(\mathrm{ftb}, \mathrm{ftc}, \mathrm{fto}, \mathrm{fta}\)).
All metrics are aggregated across runs using median and IQR (primary) and mean$\pm$std (secondary).
